% Výpočet políčka $r[4]$ zdola nahoru: vyzkoušíme všechny délky prvního
% dílu a vezmeme nejlepší výsledek. Kratší tyče už máme spočítané.
\begin{tikzpicture}[x=1cm,y=1cm,font=\small]

    \foreach \i/\val/\col in {0/{$0$}/{darkgreen!18},1/{$1$}/{darkgreen!18},
                              2/{$5$}/{darkgreen!18},3/{$8$}/{darkgreen!18},
                              4/{$?$}/{lightred!30}}{
        \draw[thick, fill=\col] ({\i*0.9},-0.35) rectangle ({\i*0.9+0.9},0.35);
        \node at ({\i*0.9+0.45},0) {\val};
        \node[font=\scriptsize, anchor=south] at ({\i*0.9+0.45},0.40) {$\i$};
    }
    \node[anchor=east, font=\small] at (-0.15,0) {$r$};

    % Kandidáti na hodnotu r[4]
    \node[font=\scriptsize, anchor=west]            at (5.15, 0.75) {$p_1+r_3=1+8=9$};
    \node[font=\scriptsize, anchor=west, darkgreen] at (5.15, 0.25)
        {$p_2+r_2=5+5=\mathbf{10}$};
    \node[font=\scriptsize, anchor=west]            at (5.15,-0.25) {$p_3+r_1=8+1=9$};
    \node[font=\scriptsize, anchor=west]            at (5.15,-0.75) {$p_4+r_0=9+0=9$};

    \draw[diredge, darkgreen, line width=1pt] (5.05,0.25) -- (4.60,0.05);

    \node[font=\scriptsize, anchor=north, align=center] at (2.25,-0.75)
        {kratší tyče už\\ máme spočítané};

\end{tikzpicture}
